\section*{Hoofdstuk \ref{ch:g2:from-seed-to-plant} --- Van zaad tot plant}

\begin{solution}{exo:g2:from-seed-to-plant:1}
Een piepklein plantje dat al begonnen is (de \omterm{def:g2:from-seed-to-plant:seed}{kiem} van de \omterm{def:g1:plants-around-us:plant}{plant}, met het
begin van een \omterm{def:g1:plants-around-us:parts}{wortel} en een scheut) en een voorraad voedsel om het te voeden.
\end{solution}

\begin{solution}{exo:g2:from-seed-to-plant:2}
\omterm{def:g2:from-seed-to-plant:germination}{Ontkiemen} is wakker worden en beginnen te groeien: het \omterm{def:g2:from-seed-to-plant:seed}{zaad} zwelt op,
zijn schil splijt, en het begint een \omterm{def:g1:plants-around-us:plant}{plant} te worden. De \omterm{def:g1:plants-around-us:parts}{wortel} komt er
als eerste uit en duikt omlaag.
\end{solution}

\begin{solution}{exo:g2:from-seed-to-plant:3}
Water en warmte. Licht staat niet op de lijst --- het \omterm{def:g2:from-seed-to-plant:seed}{zaad} \omterm{def:g2:from-seed-to-plant:germination}{ontkiemt}
vrolijk in het donker, onder de grond.
\end{solution}

\begin{solution}{exo:g2:from-seed-to-plant:4}
Het leeft van de voedselvoorraad die in het \omterm{def:g2:from-seed-to-plant:seed}{zaad} zit gepakt --- de twee
dikke helften van de boon. Die voorraad voedt het \omterm{prop:g2:from-seed-to-plant:cycle}{kiemplantje} tot zijn
groene blaadjes het licht bereiken en het zichzelf kan voeden.
\end{solution}

\begin{solution}{exo:g2:from-seed-to-plant:5}
Bijvoorbeeld: bonen, erwten, linzen, rijst, zonnebloempitten,
appelpitjes, een perzikpit.
\end{solution}

\begin{solution}{exo:g2:from-seed-to-plant:6}
Omdat een \omterm{def:g2:from-seed-to-plant:seed}{zaad} naast water ook warmte nodig heeft: in koude winteraarde
zou het blijven wachten, of rotten, in plaats van \omterm{def:g2:from-seed-to-plant:germination}{ontkiemen}. De lente
brengt de warme grond waar het \omterm{def:g2:from-seed-to-plant:seed}{zaad} op wacht.
\end{solution}

\begin{solution}{exo:g2:from-seed-to-plant:7}
\omterm{def:g2:from-seed-to-plant:seed}{Zaad}, \omterm{prop:g2:from-seed-to-plant:cycle}{kiemplantje}, bladerrijke \omterm{def:g1:plants-around-us:plant}{plant}, bloeiende \omterm{def:g1:plants-around-us:plant}{plant}, peul met nieuwe bonen.
\end{solution}

\begin{solution}{exo:g2:from-seed-to-plant:8}
Beide \omterm{def:g1:plants-around-us:parts}{wortels} groeien omlaag en beide scheuten omhoog, hoe de bonen ook
lagen. Het \omterm{prop:g2:from-seed-to-plant:cycle}{kiemplantje} vindt altijd boven en onder --- precies wat de
opmerking voorspelt.
\end{solution}

\begin{solution}{exo:g2:from-seed-to-plant:9}
In de la hadden de zaden geen water, dus konden ze niet \omterm{def:g2:from-seed-to-plant:germination}{ontkiemen}; maar
een goed \omterm{def:g2:from-seed-to-plant:seed}{zaad} gaat zonder water niet dood --- het wacht. Een jaar later
werden de wachtende zaden met water en warmte gewoon wakker.
\end{solution}

\begin{solution}{exo:g2:from-seed-to-plant:10}
De voedselvoorraad van het \omterm{def:g2:from-seed-to-plant:seed}{zaad} --- de twee dikke helften van de boon
--- speelt de rol van de dooier. \omterm{def:g2:animals-grow-up:born}{Ei} en \omterm{def:g2:from-seed-to-plant:seed}{zaad} zijn allebei het begin van
een nieuw \omterm{def:g1:living-or-not:living}{levend wezen}, volgepakt met het voedsel dat het jong nodig
heeft voordat het zichzelf kan voeden.
\end{solution}
